\documentclass[10pt,a4paper]{article}
\usepackage{fontspec,amsmath,amssymb,graphicx}
\usepackage[margin=18mm]{geometry}
\setmainfont{Times New Roman}
\setlength{\parindent}{0pt}
\setlength{\parskip}{5pt}
\allowdisplaybreaks
\sloppy
\emergencystretch=2em
\begin{document}
\section*{BT\_4\_2\_3\_01 — Lời giải bổ sung}
\subsection*{Kiến thức và trình tự giải}
\textbf{Bước 1 — Viết toán tử và biên parabol.} Đặt $P=\partial_t-a(x,t)\partial_{xx}-b(x,t)\partial_x-c(x,t)$ với $a\ge0$, $c\le0$. Biên parabol gồm mặt $t=0$ và các biên không gian; mặt $t=T$ vẫn thuộc miền được ước lượng. Với phương trình nhiệt, $a=K>0$, $b=c=0$.

\textbf{Bước 2 — Loại cực đại dương bằng dấu đạo hàm.} Tại cực đại dương nội miền, $w_x=0$, $w_{xx}\le0$, $w_t=0$; nếu cực đại trên mặt thời gian trên thì đạo hàm một phía cho $w_t\ge0$. Suy $Pw\ge0$, mâu thuẫn $Pw<0$. Với $Pw\le0$, xét $w_\varepsilon=w-\varepsilon e^{\gamma t}$, $\gamma>0$: $Pw_\varepsilon=Pw-\varepsilon(\gamma-c)e^{\gamma t}<0$. Chứng minh $w_\varepsilon\le0$ rồi cho $\varepsilon\downarrow0$.

\textbf{Bước 3 — So sánh hai hàm bằng hiệu.} Để chứng minh $v\le u\le V$, kiểm tra $Pv\le Pu\le PV$ và thứ tự của dữ liệu đầu, biên. Áp dụng bước 2 cho $v-u$ và $u-V$. Nếu có nguồn dương, hằng số không nhất thiết là nghiệm trên, nên phải tính $PV-Pu$; không dùng cực đại của dữ liệu để chặn một bài có nguồn mà chưa kiểm tra.

\textbf{Bước 4 — Xử lý Neumann, Robin và tuần hoàn.} Neumann thuần nhất có thể phản xạ chẵn qua đầu mút; điểm biên trở thành điểm trong. Robin viết theo pháp tuyến ngoài $\partial_\nu u+\sigma u=0$, $\sigma\ge0$. Tại cực đại dương không hằng, bổ đề Hopf cho $\partial_\nu u>0$, trái với $\partial_\nu u=-\sigma u\le0$. Với Robin phải thêm mức $0$ vào cận hằng nếu dữ liệu mang một dấu. Biên tuần hoàn ghép hai đầu thành đường tròn; không còn biên không gian.

\textbf{Bước 5 — Miền vô hạn và nguồn bị chặn.} Trên toàn trục/nửa trục, cần nghiệm bị chặn trên mỗi dải thời gian (hoặc điều kiện tăng trưởng đủ dùng). Dùng rào $\varepsilon((x-x_0)^2+2Kt)$ trên đoạn lớn, áp dụng cực đại trên miền hữu hạn rồi cho $\varepsilon\downarrow0$; cách này loại ảnh hưởng từ vô cực. Nếu $|Pu|\le N$ và $|u|\le M$ trên biên parabol, đặt $V=M+Nt$. Vì $PV=N-cV\ge N$, suy $|u|\le M+Nt\le M+NT$. $M$ ở đây là cận của $|u|$, không phải cực đại có dấu.
\subsection*{Lời giải áp dụng cho bài}
\textbf{Bước 1. Phát biểu cận riêng cho từng biên.} Với thời gian $[0,T]$, đặt $m_f=\min f$, $M_f=\max f$. Ở (a),
\[\min(m_f,\inf_{0\le s\le T}\mu(s))\le u\le
\max(M_f,\sup_{0\le s\le T}\mu(s)).\]
Ở (b), thêm mức $0$ do biên Robin:
\[\min(0,m_f,\inf_s\nu(s))\le u\le\max(0,M_f,\sup_s\nu(s)).\]
Ở (c),(d), $\min(0,m_f)\le u\le\max(0,M_f)$. Ở (e), $m_f\le u\le M_f$.

\textbf{Bước 2. Kiểm tra các nghiệm trên hằng.} Gọi cận trên là $C$, $w=u-C$. Ở (a), $w_x(0)=0$, $w(L)=\mu-C\le0$. Ở (b), pháp tuyến trái là $-\partial_x$, nên $(\partial_\nu+1)w=-C\le0$ tại $0$, và $w(L)=\nu-C\le0$. Ở (c), $w_x(0)=0$, $(\partial_\nu+1)w=-C\le0$ tại $L$. Ở (d), biểu thức Robin này bằng $-C\le0$ ở cả hai đầu. Lý do chọn $C\ge0$ ở Robin nằm chính ở dấu $-C$.

\textbf{Bước 3. Loại cực đại dương.} Trong miền, dùng $w-\varepsilon e^{\gamma t}$ như chứng minh cực đại yếu. Biên Dirichlet không dương. Neumann phản xạ chẵn để trở thành điểm trong, hoặc dùng Hopf; Robin ở cực đại dương không hằng có $\partial_\nu w>0$, nên $\partial_\nu w+w>0$, trái với biểu thức không dương. Trường hợp hằng dương cũng bị dữ liệu đầu và/hoặc Robin loại. Biên (e) ghép thành đường tròn vì cả giá trị và đạo hàm bằng nhau; cực đại tại đường ghép vẫn là cực đại trong. Suy $w\le0$.

\textbf{Bước 4. Cận dưới.} Áp dụng cận trên cho $-u$, với dữ liệu $-f,-\mu,-\nu$, rồi đổi dấu. Điều này trả lại đúng năm cận đã phát biểu, gồm các mức $0$ cần thiết ở Robin.\end{document}